\section{Introduction}\label{introduction}

\textbf{Problem.} We ask, for each behavior a deployed English tutor
needs, whether a sufficiently explicit system prompt is enough or the
behavior must be demonstrated in fine-tuning. A deployed tutor is
governed by a long system prompt that already \emph{states} most of what
it should do --- stay in character, ground culture in the learner's
locale, refuse role swaps, give one short example rather than a grammar
lecture, escalate to a session-end signal under sustained abuse. If
stating a behavior were sufficient, the data-side problem would be
trivial.

\textbf{Why it matters.} It is not: the behaviors split cleanly into
ones a prompt clause elicits and ones it only describes, and knowing
which is which is exactly the decision a practitioner faces when
building a task-specific model on consumer hardware --- spend the data
budget only where a prompt will not do. The tutor setting makes the
question answerable because each target behavior corresponds to a clause
already present in a realistic deployment prompt (§3): the language of
instruction, lesson topic, role structure, persona, pedagogical
contract, locale frame, and general appropriateness.

\textbf{Gap.} The prompting-versus-fine-tuning trade-off has been
studied at the level of \emph{general} alignment --- that a small
demonstration set installs broad instruction-following
\citep{zhou2023lima, ouyang2022instructgpt}, and that in-context
demonstrations often convey format more than new capability
\citep{min2022rethinking} --- but not resolved \emph{per behavior} for a
deployed task model. Synthetic-instruction pipelines and tutor-LLM
datasets (§2) target general instruction-following with a uniform
recipe; none asks, per behavior, whether the deployment prompt already
elicits what the data teaches.

\textbf{Approach and contributions.} We supply the same fully-specified
deployment prompt at evaluation to a fine-tuned 0.8B student and to a
ladder of prompt-only baselines (0.8B base, 0.8B instruct, 4B instruct,
and the 9B teacher that generated the training data;
Figure\textasciitilde{}\ref{fig:design}). Because the instruction is
present for every condition, a prompt-only failure cannot be blamed on
under-specification. Our contributions:

\begin{enumerate}
\def\labelenumi{\arabic{enumi}.}
\tightlist
\item
  \textbf{A per-capability map of the train-versus-prompt boundary}
  (Figure\textasciitilde{}\ref{fig:boundary}): promptable when a single
  clause both \emph{describes} and \emph{elicits} the behavior (locale
  fidelity, role-swap and topic re-anchoring), not promptable when it
  requires cross-turn state-tracking (persistence) or the suppression of
  a strong competing prior (pedagogical withholding).
\item
  \textbf{An evidenced negative result about the conventional metric}:
  per-axis redirect F1 is a context-blind \emph{type} classifier that
  ties a 0.8B student with the 9B teacher; we replace it with matched
  per-capability instruments (§5).
\item
  \textbf{Reusable apparatus}: a locale-aware, yield-aware generation
  pipeline (§3) and a \emph{partially validated} trigger-position
  decorrelation construction for the persistence sentinel (§5).
\end{enumerate}

All experiments run on one RTX 3060 12GB GPU --- the regime where the
boundary matters most, since a large general-purpose model is not
deployable there. We evaluate one family (Qwen; a genuine limitation,
§6) and one locale (which does not limit the persistence/withholding
boundary, both structural behaviors independent of the locale backdrop).
